% 本文件由 LaTeX 排版 Agent 根据有效 translation.json 自动生成。
% author=Gregory of Nyssa; work_id=2910; title=论基督的圣洗

\chapter{论基督的圣洗}

% structure: p0001 source=h1 target=omitted reason=page-title-already-rendered-by-parent

一篇为\OriginalTerm{光明节}{Lights}而作的讲道。\par

现在我认出我的羊群：今日我看见教会的惯常面貌，当你们厌恶地回避甚至肉身的挂虑，以不减的人数聚集为天主服务——当人群挤满房屋，进入神圣的圣所，当找不到位置的众人像蜜蜂一样填满外面的空地。因为他们中有些在屋内工作，有些在外面绕着蜂巢嗡嗡作响。要如此作，我的孩子们：永远不要放弃这种热忱。因为我承认我感受到牧者的情感，当我被安置在这个瞭望塔上时，我希望看见羊群聚集在山脚周围；当我得见此景时，我便充满奇妙的热情，愉快地投入讲道，如同牧者吹奏乡村曲调一般。但若情况相反，你们在远方的游荡中迷途，就像你们刚在上一个主日所作的那样，我便十分忧愁，乐于沉默；我便考虑从此处逃离，寻求厄里亚先知的加尔默耳，或某个无人居住的岩石；因为抑郁之人自然选择孤独和独处。但如今，当我看见你们带着全家蜂拥而至时，我便想起先知的话，那是依撒意亚从远处宣告的，预先对教会——这位美丽且儿女众多的母亲——所说的话：\Quote{这些如云彩飞来、又如鸽子衔雏向我飞来的是谁？}是的，他还加上这句话：\Quote{地方太狭窄，请给我地方居住\BibleRef{依 49:20}}。因为这些预言，圣神的能力指向了天主那人口众多的教会，它后来要充满整个大地，从地极到地极。\par

这时辰已经来到，它在进程中带来了对神圣奥迹的纪念，这些奥迹洁净人——这些奥迹从灵魂和肉体中清除那难以洗净的罪，并将我们带回我们原初的美善，那是天主——至高的工匠——铭刻在我们身上的。因此，你们这些已受洗的人聚集在一起；你们也带来了那些尚未体验过的人，如同良善的父亲，以细心的引导带领未受洗的人达至信德的完全领受。至于我，我为两者都欢喜——为你们已受洗的，因为你们被丰富的恩赐所充实；为你们未受洗的，因为你们拥有希望的佳美期待——罪债的赦免，束缚的解除，与天主的亲密关系，自由的宣讲之胆量，以及取代奴役的与天使同等的地位。因为这一切，以及由此而来的一切，圣洗圣事的恩宠为我们确保并传递给我们。因此，让我们将圣经的其他问题留给其他场合，坚守我们眼前的主题，尽我们所能奉献与这庆节相称的礼物：因为每个庆节都有其特定的处理方式。所以，我们用婚宴的歌迎接婚礼；对哀悼，我们带来与葬礼相称的奉献；在事务时期，我们严肃地说话；在喜庆时期，我们放松精神的专注和紧张；但每个时期我们都保持不受与其性质相异之事的干扰。\par

基督，说祂是几天前出生的——祂的生出是在一切可感和可理解的事物之前。今天祂由若翰受洗，为要洁净那被玷污的，为要从高天带来圣神，并将人高举到天上，使那堕落者被举起，那引人跌倒者蒙羞。不要惊讶天主在我们的事上显示了如此大的热忱：因为那亏负我们的人精心策划了陷害我们的阴谋；而我们的创造主则以先见之明拯救了我们。那邪恶的诱惑者，针对我们人类编造了新的罪计，拖着他的蛇队，一个与他意图相称的伪装，带着他的不洁进入那与他相似的东西——那爬行之物，他在意志上是属地和属世的。但基督，祂恶行的修复者，取了完整的人性，拯救了人，成为我们众人的预像和模式，以圣化每一行动的首果，并使祂的仆人在热诚遵行传统上毫无疑虑。那么，圣洗圣事是罪的净化，过犯的赦免，更新和重生的原因。所谓重生，是指理智中构想的重生，而非肉眼所能辨识的。因为我们不会像犹太人尼苛德摩及其有些迟钝的理解那样，将老人变成孩童，也不会将满面皱纹、头发花白的人重塑为柔嫩年少，倘若再把人带回母腹中；但我们确实借君王般的恩宠，将那带着罪痕、在恶习中衰老的人带回到婴儿的无辜。因为如同初生婴儿不受控告和刑罚一样，重生之子也无需负责，因君王般的宽厚而免于责任。这恩赐并非水所给予（因为若如此，水就成为高于万物者），而是天主的命令，以及圣神以圣事方式降临而释放我们。但水用以表达洁净。因为我们习惯用水洗涤，使身体在沾染灰尘或泥土后得洁净，因此在圣事行动中也使用水，并藉我们感官所及之物显示属灵的光辉。然而，若认为合适，让我们继续更详尽、更深入地探究圣洗圣事，从圣经的宣告出发，如同从泉源开始：\BibleQuote{除非人由水和圣神而生，不能进天主的国。}为什么二者都被提及？为什么单凭圣神不被认为足以完成圣洗？我们知道得很清楚，人是复合而非单一的：因此，对那双重且混合的人，指定了同源且相似的药物来医治——对他可见的身体，水，可感的元素；对他不可见的灵魂，不可见的圣神，由信德呼求，以不可言喻的方式临在。因为\BibleQuote{风随意向哪里吹，你听到它的响声，却不知道它从哪里来，往哪里去。}祂祝福受洗的身体和施洗的水。因此，不要轻视这神妙的圣洗池，也不要因它用了水而将其视为平常。因为运行的大能是强大的，由此成就的事是奇妙的。因为我现在在这祭台旁边，它本质上是普通的石头，与建造房屋、装饰地面的其他石板毫无区别；但它被祝圣以服务于天主，并接受了祝福，就成了神圣的供桌，不受玷污的祭台，不再被众人手触，而只由司铎以恭敬之心触摸。同样，饼起初是普通的饼，但当圣事行动祝圣它时，它就被称为并成为基督的身体。圣事之油也是如此；酒也是如此：在祝福之前它们价值微小，但在圣神赋予圣化之后，每一种都有其独特的功效。同样，圣言的德能也使司铎成为可敬和尊贵的，藉新赋予他的祝福，与大众的共同体分离。虽昨日他还是大众之一、百姓之一，却突然被造就为引导者、主持者、正义的教师、隐秘奥秘的教导者；他如此行，身体或外貌毫无改变；但尽管在外表上仍是原先那个人，却藉某种不可见的大能和恩宠，就他不可见的灵魂而言被转化为更高的状态。因此有许多事物，你若细察，会发现它们的外表卑贱，但所成就的却伟大；尤其当你从古史中搜集与我们探讨的主题同源相似的实例时。\Person{梅瑟}的棍杖是一根榛树枝。那不过是普通的木头，人人都可砍伐、搬运、随意加工，并任意投入火中。但当天主乐意藉那棍杖成就那些崇高、言语无法表达的奇事时，那木头就变成了蛇。又有一时，他击打水，就使水变为血，又使无数的青蛙涌出；他又分开海洋，从深处劈开，水不再合拢。同样，一位先知的外衣，虽不过是山羊皮，却使\Person{厄里叟}名扬天下。又十字架的木头对众人有救恩之效，虽据我所知，它是一棵劣树的木头，比不上大多数树木的价值。所以，荆棘丛向\Person{梅瑟}显示天主临在的显现；\Person{厄里叟}的遗骨使死人复活；泥土使胎生瞎子复明。所有这些事物，虽是无灵魂、无感觉的物质，但在接受天主的大能后，就成了成就伟大奇迹的工具。如今用类似的推理，水也一样，虽不过是水，但当自上而来的恩宠祝圣它时，它使人更新，获得属灵的重生。若有人又提出难题，带着疑问和怀疑，不断追问水及其中所行的圣事行动如何使人重生，我很公正地回答他：\Quote{请给我指出肉身生育的方式，我就向你解释灵魂中重生的能力。}你或许会说，作为解释：\Quote{是种子使人产生。}那么反过来从我们学习：祝圣过的水洁净并光照人。若你又用你的\Quote{如何？}反驳我，我将更激烈地回应：\Quote{那稀稀的、无形的质料如何成为人？}如此，这辩论在推进中会涉及整个受造界的每一事物。天如何存在？地如何？海如何？每一事物如何？因为处处人的理智在寻求发现中困惑，便退回到这个音节\Quote{如何}，如同不能行走的人退回到座位上。简言之，处处天主的大能和作为都是不可理解、不能被归结为规则的，祂轻易地产生祂所愿意的一切，同时向我们隐藏祂作为的细节知识。因此，有福的\Person{达味}将心智集中于创造的宏伟，内心充满困惑的惊叹，说出那首众人咏唱的诗歌：\BibleQuote{上主，祢的化工何其繁多！祢以智慧成就了一切。}他感知了智慧，却未能发现智慧的艺术。那么，让我们放下探究超出人类能力之事的任务，而去寻求那些似乎部分在我们理解范围内的事：为什么洁净是由水来实现的？三次浸入的目的是什么？教父们所教导、我们心智所接纳同意的如下：我们承认构成世界的四元素，这是人人都知道的，即便不说出它们的名字；但若为了更简单起见，告诉你们它们的名字，那就是火、气、土、水。现在，我们的天主和救主，为了我们而完成救恩计划，进入了这四元素中的第四个——土，以便从那里使生命升起。我们领受圣洗，效法我们的主、师傅和引导者，并非被埋葬在土里（因为那是完全死亡的躯体的庇护所，覆盖我们本性的软弱和腐朽），而是来到与土相似的元素——水中，将自己隐藏在其中，如同救主在土中一样；这样做三次，我们就为自己代表了那在三天内完成的复活恩宠；我们这样做并非在沉默中领受圣事，而是在我们身上诵念我们信仰、也盼望的三位神圣位格的名字，从祂们而来我们现世和未来的存在。你，那大胆对抗圣神荣耀的人，也许感到冒犯，并且吝啬给予圣神那由虔诚者所献的敬拜。停止与我争辩：你若能，就反对那些给予人类圣洗呼求规则的主的话。主的命令说什么？\BibleQuote{你们要去使万民成为门徒，因父及子及圣神之名给他们授洗。}\BibleRef{玛 28:19}因父之名如何？因为祂是万物的原始原因。因子之名如何？因为祂是创造的创造者。因圣神之名如何？因为祂是成全一切的德能。因此我们向圣父俯首，为能成圣；也向圣子俯首，为达到同一目的；也向圣神俯首，为使我们事实上和名义上成为祂所是的。圣化中没有区别，即圣父圣化更多、圣子较少、圣神比另外两位更少。那么，你为什么将三位分割成不同本质的碎片，制造三个彼此不同的神，而你从三者都领受同一恩宠？\par

然而，例证总能使论述对听者更为生动，故我拟借一比喻来教导这些亵渎者的心智，藉由属地和卑微的事物，解释那些伟大且感官不可见的事。若你竟遭遇在仇敌中为囚徒的不幸，身处桎梏与苦难，为那昔日曾有的自由而叹息——且突然有三位男士，他们是暴虐主子国中的显要公民，共同解除你身上的束缚，平均付出你的赎金，并将赎款分摊于他们之间——那么，你蒙此恩惠，岂不将三人同等视为恩人，并按均等份额偿还赎金，因为他们为你所受的劳苦与花费是共同的——倘若你是公正评价这恩惠的话？按比喻而言，我们可看出这一点，因我们目前的目的不是严格陈述信仰。让我们回到当前的节期，以及它呈现给我们的主题。\par

我发现，不仅是在十字架受难后写成的福音宣扬圣洗的恩宠，而且，甚至在我们主的降生成人之前，古老的圣经就已处处预表了我们重生的样式；并非清楚地显明其形式，而是以隐晦的言论预示了天主对人的爱。就如羔羊曾被预告，十字架也曾被预告，同样，圣洗也藉行动和言语被预示。让我们向喜爱善思的人回忆这些预象——因为节期之故，必须记起它们。\par

亚巴郎的婢女哈加尔——保禄在训导迦拉达人时曾以寓意解释——因撒辣的恼怒被主母逐出家门（因为婢女被怀疑与主人有关，是正室所难以忍受的），她带着怀中婴儿依市玛耳流落到荒凉之地。当生活所需匮乏，自己濒临死亡，孩子更为危急时，因皮囊里的水用尽了（不可能让曾居于永泉表像中的会堂拥有维持生命所需的一切），一位天使忽现，指给她一眼活水井，她从中汲水，救了依市玛耳。看，这是一个圣事性的预象：从起初，正因活水，那垂死之人得救——这水原不存在，而是藉天使赐予的恩赐。后来，那曾使依市玛耳与母亲被逐出父家的依撒格将要成婚。亚巴郎的仆人去为婚事说媒，为主人寻得新娘，在井边找到了黎贝加；这桩将产生基督后裔的婚姻，其开始和初次盟约都在水中。是的，依撒格本人牧放羊群时，也在旷野各处掘井，却被异族人堵塞填满，这预表了后来所有阻碍圣洗恩宠、在反抗真理时大声喧哗的不虔敬之人。然而殉道者与司铎们通过掘井战胜了他们，圣洗的恩赐溢满普世。按经文同样之意，雅各伯也匆忙寻妻，在井边无意中遇到辣黑耳。井口盖着一块大石，许多牧人常聚来一起滚开石头，然后饮自己和羊群。但雅各伯独自滚开石头，饮了新娘的羊群。我想，这事是隐晦之语，是未来之事的影儿。因为那盖着的石头，不就是基督自己吗？依撒意亚论及他说：\BibleQuote{我要在熙雍的地基上放一块宝石，一块精选的、宝贵的宝石，}达尼尔也说了：\BibleQuote{一块石头未经人手而被凿出来，}即基督没有男人而诞生。因为石头从岩石中被凿出，无需凿石者或石匠，这是新奇惊人的事；同样，未婚童贞女生子，是超乎一切奇迹的事。所以，井上盖着灵磐石基督，他将重生之洗藏在深处和奥秘中，需要很长时间——如一根长绳——才能显露。除了以色列（即看见天主的理智），无人能滚开这石头。但他既汲水，又给辣黑耳的羊群饮；即他揭示了隐藏的奥秘，将活水赐给教会的羊群。此外，还要加上雅各伯三根木杖的历史。因为自从三根木杖放在井旁，多神教徒拉班从此变穷，而雅各伯变得富有，牲畜成群。现在，让拉班被解释为魔鬼，雅各伯为基督。因为圣洗建立后，基督夺走了撒殚所有的羊群，自己变得富足。再者，伟大的梅瑟，当他还是俊美婴儿，尚在襁褓时，因法郎冷酷残忍对男婴下的普遍残酷法令而遭难，被放在河岸边——并非赤身，而是放在一个箱子里（因为法律理应被封闭在柜中预表）。他被放在水边；因为法律和希伯来人每日的洒水礼，稍后在完美而奇妙的圣洗中得到显明，它们离恩宠很近。又，按受默感的保禄的看法，百姓本身经过红海，宣告了藉水得救的福音。百姓过去，而埃及王及其军队被吞没，藉这些行动预表了这圣事。因为即使如今，每当百姓进入重生之水，逃离埃及、脱离罪恶的重担时，便得释放而得救；但魔鬼及其仆役（我指恶灵）则因悲伤而窒息，灭亡了，视人的得救为自己的不幸。\par

即便这些例子已足够肯定我们目前的立场；但爱好善思的人仍不可忽略接下来的内容。据我们所知，希伯来人在经历许多苦难，并在旷野中完成他们疲惫的旅程之后，直到在\Person{若苏厄}的引导和生命导航之下，被带到\Place{约但河}的渡口，才进入应许之地。但显然，在那溪流中竖立十二块石头的\Person{若苏厄}，也预示了十二位宗徒——圣洗圣事的施行者——的来临。再者，那位年迈的提市贝特人的奇妙祭献（超出人类的理解），除了以行动预表对圣父、圣子、圣神的信德和救赎之外，还能做什么呢？因为当全体希伯来人践踏了他们祖先的宗教，陷入\OriginalTerm{多神论}{polytheism}的错误，他们的国王\Person{阿哈布}被\OriginalTerm{偶像崇拜}{idolatry}所迷惑，与那名字不祥的\Person{依则贝耳}（他邪恶的生活伴侣和亵渎之心的卑鄙煽动者）为伍时，那位充满圣神恩宠的先知，来与阿哈布会面，在国王和全体人民面前，与巴耳的神职人员进行了一场奇妙而惊人的较量。他向他们提出不用火祭献公牛的任务，使他们显出可笑而凄惨的窘境，徒然地向不存在的神祇祈祷和呼喊。最后，他自己呼求他自己的、也是真天主的那位，以进一步的夸张和增添完成了所提出的考验。因为他并非仅仅藉祈祷使火从天降在干柴上，而是敦促并命令侍从带来大量的水。当他将桶中的水三次倾倒在那劈开的柴上时，藉着祈祷他从水中点燃了火，为的是通过元素的对立（却又如此友好地合作）来以超绝的力量彰显他自己的天主的能力。现在，藉着那奇妙的祭献，\Person{厄里亚}清晰地为我们预示了日后将设立的圣洗圣事的礼仪。因为火被三次倒上的水所点燃，所以清楚表明，哪里有这样的\OriginalTerm{神秘之水}{mystic water}，那里就有那燃点、温暖而炽热的圣神，祂焚烧不敬虔的人，光照信友。是的，还有他的门徒\Person{厄里叟}，当患麻疯病的叙利亚人\Person{纳阿曼}来向他恳求时，他在\Place{约但河}中为他洗涤而洁净了病人，清楚指示了将要来临的事，既通过普遍用水的做法，也特别通过在河中的浸洗。因为\Place{约但河}是唯一一条河流，在自己内接受了圣化和祝福的初果，通过它的水道将圣洗的恩宠传遍全世界，仿佛从某源泉中通过它自己提供的预像。这些就是藉圣洗而重生的行为和实践中的预兆。其余的，让我们考虑它在言语和语言中的预言。\Person{依撒意亚}呼喊道：\BibleQuote{你们要洗濯，洁净自己，从你们的灵魂中除去邪恶。}\Person{达味}说：\BibleQuote{你们要亲近祂，被光照，你们的脸就不会羞愧。}而\Person{厄则克耳}，写得比两者都更清晰明白，说：\BibleQuote{我要在你们身上洒清水，你们就洁净：我要洗净你们脱离各种不洁和各种偶像。我还要赐给你们一颗新心，在你们五内注入一种新精神：我要从你们的肉体内除去铁石的心，给你们换上一颗血肉的心，并将我的神放在你们内。}\Person{匝加利亚}也非常明显地预言了关于\Person{若苏厄}的事，他穿着污秽的衣服（即仆人的肉身，也就是我们的肉身），并脱去他不体面的衣裳，为他穿上干净美丽的衣服；藉着这比喻的说明教导我们：在耶稣的圣洗中，我们所有人都脱去罪恶如同破烂的补丁衣服，穿上了神圣而最美丽的重生之衣。我们又将\Person{依撒意亚}的那句神谕置于何处？他向旷野呼喊：\BibleQuote{干渴的旷野，喜乐吧；沙漠要欢欣，如百合花盛开；约但的荒芜之地要开花，要喜乐。}显然，他并非向没有灵魂或感觉的地方宣告喜乐的好消息：而是藉着沙漠的比喻，向那干渴、毫无装饰的灵魂说话，正如\Person{达味}所说：\BibleQuote{我的灵魂在祢面前像干涸的土地}和\BibleQuote{我的灵魂渴慕大能者，生活的天主。}同样，主在福音中说：\BibleQuote{谁若渴，让他到我这里来喝。}又对撒玛黎雅妇人说：\BibleQuote{凡喝这水的，还要再渴；但谁若喝了我赐给他的水，将永远不渴。}\BibleRef{若 4:13}而\BibleQuote{加尔默耳的光辉}\BibleRef{依 35:2}赐给那具有旷野模样的灵魂，即藉圣神赐予的恩宠。因为\Person{厄里亚}曾住在\Place{加尔默耳}，那山因住在那里的圣德而成名；此外，\Person{若翰洗者}——充满厄里亚的精神——圣化了\Place{约但河}；因此先知预言\BibleQuote{加尔默耳的光辉}应赐给那河。并且，他将\BibleQuote{黎巴嫩的荣耀}\BibleRef{依 35:2}藉其高大树木的相似比喻转移到那河。因为正如伟大的\Place{黎巴嫩}在其所产生并滋养的树木本身呈现了充分的惊奇理由，同样，\Place{约但河}因使人重生并将他们栽种在天主的乐园中而得荣耀；关于他们，如同圣咏作者的话所说，常青并结满德行之叶，\BibleQuote{叶子不会枯萎}，天主将欢喜，按时收他们的果实，如同一个好农夫，为自己的工作而喜乐。受感动的\Person{达味}，也预言了圣父在圣子受洗时从天上发出的声音，为引导那些直到那时只凭感官看到祂人性的卑微状态的听众，达至属于天主性体本身的尊贵，在他的书中写下了这段经文：\BibleQuote{上主的声音响彻在水上，上主的声音威震八方。}但在此我们须结束来自神圣圣经的见证：因为若有人试图详细挑选每段经文并将其在一本书中展示，论述将无限延长。\par

但是你们所有因重生的恩赐而喜乐，并以那救恩的更新而夸耀的人，都要在领受圣事的恩宠之后，向我展示你们行为上应有的改变，并以你们纯洁的言行，表明你们因这向善的转变所发生的不同。因为在我们眼前的事物中，没有一样改变：身体的特征保持不变，可见本性的形态也毫无不同。但是确实需要某种明显的证据，藉以认出那新生的人，通过清晰的标记分辨新人与旧人。我认为这些标记可以在灵魂的有意动向中找到，灵魂藉此脱离它旧有的习惯生活，进入一种新的行为方式，并将清楚地教导那些认识它的人，告诉他们已经变得与从前不同，身上不再带有任何识别旧我的标记。这，如果你们听从我的劝告，并将我的话视为法则，便是转变的方式。那在领受圣洗之前的人，曾是放纵的、贪婪的、攫取他人财物的、辱骂人的、说谎的、诽谤的，以及一切与这些相关并从这些衍生的事。现在让他成为有秩序的、清醒的、满足于自己所有，并从其中分施给贫穷的人，诚实的、谦和的、亲切的——总之，追随一切可赞美的行为方式。因为如同黑暗被光明驱散，黑色被白色覆盖而消失，同样，那旧人也在披戴义德的工作时消失。你们看\Person{匝凯}，也藉着生活的改变杀死了税吏，以四倍偿还那些被他冤枉损害的人，并将余下的分给穷人——那些他从前藉不法手段从他所压迫的穷人那里获得的财富。福音作者\Person{玛窦}，另一位税吏，与\Person{匝凯}从事同样的行业，在蒙召之后立刻改变了他的生活，如同摘下面具一般。\Person{保禄}曾是迫害者，但在领受恩宠后成为宗徒，为了基督的缘故忍受锁链的重负，以此作为补赎和悔改的行动，为他从前从法律领受并用来反对福音的那些不义的枷锁。你们在重生中应当如此：你们应当如此消除那些倾向罪恶的习惯；天主的子女应当如此生活：因为在领受恩宠后，我们被称为祂的儿女。因此我们应该仔细审视我们圣父的特性，藉着将自己塑造并形成与圣父相似的模样，我们才能显为那按恩宠召叫我们获得义子身份者的真正儿女。因为私生子和冒名顶替的儿子，其行为违背了他父亲的高贵，是可悲的耻辱。因此，我以为主自己也在福音中为我们制定了生活法则，对祂的门徒说了这些话：\BibleQuote{你们要爱你们的仇敌，为那迫害你们的人祈祷；这样你们就可以成为你们在天上圣父的儿女；因为祂使太阳上升，光照恶人也光照善人；降雨给义人，也给不义的人。}\par

因此，在义子身份的尊荣之后，魔鬼也越发猛烈地谋害我们，它怀着嫉妒的目光，看到那新生之人的美丽，热切地趋向那它曾坠落的天上之城，便消瘦憔悴；它向我们燃起炽热的诱惑，竭力想要剥夺我们的第二层装饰，正如它剥夺我们最初的装束一样。但是当我们意识到它的攻击时，我们应该对自己重复宗徒的话：\BibleQuote{我们这许多受过洗归于基督耶稣的人，就是受洗归于祂的死亡 \BibleRef{罗 6:3}。}如今，如果我们已与祂的死亡同化，那么罪恶在我们内确实成了一具死尸，被圣洗的长矛刺穿，如同那个行淫的人被热心的\Person{丕乃哈斯}刺穿一样。\BibleRef{户 25:7}因此，不祥者，从我们这里逃走吧！因为你想要掠夺的是一具死尸，一具早已与你有染、早已为快乐而丧失理智的死尸。死尸不会迷恋肉体，死尸不会被财富迷惑，死尸不诽谤，死尸不说谎，不夺取不属于自己的东西，不辱骂遇到的人。我的生活方式是为另一种生命而规范的：我已学会轻视世上的事物，经过地上的事物，赶向天上的事物，正如\Person{保禄}明确作证说，世界已为他钉在十字架上，他也为世界钉在十字架上。这些是真正重生之灵魂的话语：这些是新受洗之人的心声，他记得自己在领受圣事时向天主所做的宣认，许诺为了对祂的爱而轻视所有的痛苦和快乐。\par

如今，关于今日的圣题，我们已经说得足够多了；这圣题是循环的岁月在指定的时期带给我们的。我们将善用余下的时间，把我们的讲论转向那慈爱的施予者，祂赐予如此宏恩；我们献上几句言语，作为对大事的回报。因为主啊，祢确实是纯善的永恒泉源，祢曾公义地转面离弃我们，却又以慈爱怜悯了我们。祢曾憎恨，却又和好；祢曾诅咒，却又祝福；祢将我们逐出乐园，又召回我们；祢剥去那无花果树叶的不体面遮盖，给我们穿上贵重的衣裳；祢打开监牢，释放了被定罪的人；祢用清水洒我们，洗净我们的污秽。\Person{亚当}不再因祢的呼唤而羞惭，也不再因良心谴责而躲藏，蜷缩在乐园的丛林中。那旋转的火焰之剑也不再环绕乐园，使走近者无法进入；相反，一切对于我们这些罪的继承者，都化为喜乐：乐园，甚至天堂本身，都可由人践踏；受造物，包括世界内和世界上的，曾彼此不和的，在友谊中结合；而我们人类被允许加入天使的歌咏，向天主献上赞美的敬拜。为此，让我们向天主唱那首喜乐之歌，就是昔日被圣神触动嘴唇的人高声唱出的：\BibleQuote{愿我的灵魂因主而欢乐，因为祂给我穿上了救恩之衣，给我披上了喜乐之袍；好像新郎，祂给我戴上冠冕；又好像新娘，祂以华服装饰我。}诚然，装饰新娘的那位就是\Person{基督}，祂昔在、今在、将来永在，从今世直到永远，当受赞美。阿们。\par

\SourceRecord{Translated by H.A. Wilson.；Nicene and Post-Nicene Fathers, Second Series；Vol. 5.；Edited by Philip Schaff and Henry Wace.；Buffalo, NY: Christian Literature Publishing Co.,；1893.；<http://www.newadvent.org/fathers/2910.htm>.}
